% !TeX root = ../main.tex
%%% ---------------------------------------------------------------------------
%%% 工作一：跨层注意力门控
%%%
%%% 目标页数：4~6 页（20 分钟答辩里，两块工作合计应占一半以上时间）。
%%% 每块工作的内部结构固定为：想法 → 方法 → 结果 → 这块工作证明了什么。
%%%
%%% 公式只放论点必需的那一两个，推导进备用页。委员会真想看推导会问，
%%% 那时候翻备用页比事先塞在正文里更显得有准备。
%%% ---------------------------------------------------------------------------
\section{工作一：跨层注意力门控}

\begin{frame}{让每个空间位置自己决定信任哪一层特征}
  \begin{block}{门控权重}
    \vspace{-0.6em}
    \begin{equation*}
      \mat{G}_l = \sigma\!\left(
        \mat{W}_2 \relu\!\left(
          \mat{W}_1 \left[ \operatorname{GAP}(\featmap_l) \,\Vert\,
                           \operatorname{GAP}(\tilde{\featmap}_{l+1}) \right]
        \right)\right)
    \end{equation*}
  \end{block}
  \vspace{0.4em}
  \begin{block}{融合}
    \vspace{-0.6em}
    \begin{equation*}
      \featmap_l' = \alert{\mat{G}_l} \odot \featmap_l
                  + (\mathbf{1} - \alert{\mat{G}_l}) \odot
                    \mathcal{D}(\tilde{\featmap}_{l+1})
    \end{equation*}
  \end{block}
  \takeaway{$\mat{G}_l \to \mathbf{0}$ 时退化为标准 FPN，所以这是它的严格推广}
\end{frame}

\begin{frame}{方差不增保证了训练稳定性}
  \begin{theorem}[方差不增]
    设两路特征互不相关、方差分别为 $\sigma_1^2,\sigma_2^2$，则对任意
    $\mat{G}_l \in [0,1]^{H_l\times W_l}$，
    \begin{equation*}
      \Var\!\left[\featmap_l'\right] \le \max\setof{\sigma_1^2,\sigma_2^2}.
    \end{equation*}
  \end{theorem}
  \vspace{0.6em}
  \begin{itemize}
    \item 逐点展开是 $g$ 的二次函数，最大值只能在端点取得
    \item 完整证明与紧性讨论见备用页
  \end{itemize}
  \takeaway{不放大方差，因此不需要额外的训练技巧来稳住}

  \note{委员会最可能在这里问「独立性假设成立吗」。
        答案：不严格成立，但实验未观察到发散；备用页有相关性的实测。}
\end{frame}

\begin{frame}{门控单独带来 2.3 个百分点，增益集中在小目标}
  \begin{columns}[T, onlytextwidth]
    \begin{column}{0.5\textwidth}
      \begin{table}
        \centering\small
        \begin{tabular}{l S[table-format=2.1] S[table-format=2.1]}
          \toprule
          {方法} & {mAP / \%} & {AP\textsubscript{S} / \%} \\
          \midrule
          FPN 基线     & 74.7 & 48.5 \\
          + 门控       & \bfseries 77.0 & \bfseries 52.6 \\
          \bottomrule
        \end{tabular}
      \end{table}
    \end{column}
    \begin{column}{0.46\textwidth}
      \centering
      \placeholder{\linewidth}{3.6cm}{figures/gate-heatmap.pdf}
    \end{column}
  \end{columns}
  \takeaway{小目标增益（+4.1）远高于总体（+2.3），与设计意图一致}
\end{frame}
